\documentclass[11pt]{article}
\usepackage[margin=1in]{geometry}
\usepackage{amsmath,amssymb,amsthm}
\usepackage[colorlinks=true,linkcolor=blue,citecolor=blue,urlcolor=blue]{hyperref}

\newtheorem*{claim}{Claim}

\begin{document}

\begin{center}
{\Large Literature-implied answer for a real \(p=3/2\) subcase}\\[4pt]
{\large arXiv:2012.02016 Question 5 via arXiv:2207.07938}
\end{center}

\paragraph{Status.}
\texttt{literature\_implied\_answer (partial real-scalar variant)}. This is not a new solution of the invariant-subspace problem and not a solution of the complex \(1<p<2\) case.

\paragraph{Source question.}
Grivaux--Matheron--Menet, \emph{Does a typical \(\ell_p\)-space contraction have a non-trivial invariant subspace?}, arXiv:2012.02016, Section 7, Question 5, asks:
\begin{align*}
&X=\ell_p,\ 1<p<2,\ \text{or }X=c_0: \\
&\qquad\text{for every fixed }M>1,\text{ is }(MT)^*\text{ hypercyclic for typical }\\
&\qquad T\in(\mathcal B_1(X),\mathrm{SOT})?
\end{align*}
For the \(\ell_p\) case it also asks whether, at least, the typical SOT contraction has no eigenvalue.

\paragraph{Supporting theorem.}
Grivaux--Matheron--Menet, \emph{Generic properties of \(l_p\)-contractions and similar operator topologies}, arXiv:2207.07938, Theorem C (restated as Theorem \texttt{toutcapourca}) proves that, for real \(\ell_p\) with \(p=3\) or \(p=3/2\), all topologies on \(\mathcal B_1(\ell_p)\) lying between WOT and SOT* are similar. In particular, SOT and SOT* have the same comeager subsets on real \(\ell_{3/2}\).

\begin{claim}
Let \(X\) be real \(\ell_{3/2}\). For every fixed \(M>1\), a typical \(T\in(\mathcal B_1(X),\mathrm{SOT})\) is such that \((MT)^*\) is hypercyclic. Consequently, a typical SOT contraction on real \(\ell_{3/2}\) has no eigenvalue.
\end{claim}

\begin{proof}
For \(1<p<\infty\), arXiv:2012.02016 records in Proposition \texttt{monsteraide} that, for every \(M>1\),
\[
\mathcal H_M=\{T\in\mathcal B_1(\ell_p): (MT)^*\text{ is hypercyclic and }\|T\|=1\}
\]
is dense \(G_\delta\) in \((\mathcal B_1(\ell_p),\mathrm{SOT}^*)\). The proof is the standard SOT/SOT* hypercyclic-density proof and is field-insensitive, so it applies to the real scalar space as well.

Now take \(X=\ell_{3/2}\) over \(\mathbb R\). By arXiv:2207.07938, Theorem C, SOT and SOT* are similar on \(\mathcal B_1(X)\); hence they have the same comeager sets. Therefore \(\mathcal H_M\), already SOT*-comeager, is also SOT-comeager.

Finally, if \(S=(MT)^*\) is hypercyclic, then \(S^*=MT\) has no eigenvalue. Thus \(T\) has no eigenvalue.
\end{proof}

\paragraph{Scope.}
This answers only the real \(p=3/2\) scalar variant/subcase of the hypercyclicity part of Question 5. It does not answer the original complex \(1<p<2\) problem, the \(c_0\) case, the question whether SOT* comeagerness transfers to SOT for all \(1<p<2\), or the main invariant-subspace question.

\paragraph{Search evidence.}
The run indexes were searched for arXiv:2012.02016, the title keywords, SOT/SOT*, hypercyclicity, and invariant-subspace phrases. The only close prior durable hit was the attempt note for arXiv:2207.07938. Web searches on 2026-07-03 for the exact 2012 title and the \(p<2\) hypercyclic/no-eigenvalue phrases found no newer complex-case resolution.

\paragraph{References.}
\begin{itemize}
\item S. Grivaux, E. Matheron, Q. Menet, \emph{Does a typical \(\ell_p\)-space contraction have a non-trivial invariant subspace?}, arXiv:2012.02016.
\item S. Grivaux, E. Matheron, Q. Menet, \emph{Generic properties of \(l_p\)-contractions and similar operator topologies}, arXiv:2207.07938.
\end{itemize}

\end{document}
